% GENERATED FILE. Edit canonical lessons, then run npm run generate:books.
% canonical-source-hash: fnv1a64:062c02c3c5e787b5
% canonical-lessons: RU-C487-syp, RU-C487-zud, RU-C487-chesatsya, RU-C487-privivka, RU-C487-zarazhatsya

\chapter{Rashes and Jabs}
\label{ch:ru-rashes-and-jabs}
\hlchaptermodality{487}

\begin{chapteropening}
\textbf{By the end of this chapter:} \emph{I can talk about rashes, itching and catching things.}

Complete the last lesson of chapter 487: I can talk about rashes, itching and catching things.
\end{chapteropening}

\section[\texorpdfstring{\ru{сыпь} (syp)}{сыпь (syp)}]{\ru{сыпь} (syp) — a rash}

\label{lesson:RU-C487-syp}



\begin{quote}
Before the new one: say the Russian for to go hoarse, then the Russian for to tickle.
\end{quote}

\subsection*{You'll want to know: \ru{сыпь}}
\textbf{\ru{сыпь}} (\emph{syp}) — \textquotedblleft{}a rash\textquotedblright{}.

\begin{quote}
\ru{После} \ru{клубники} \ru{у} \ru{ребёнка} \ru{появилась} \ru{сыпь}.
\end{quote}

\begin{quote}
\emph{Pósle klubníki u rebyónka poyavílas syp.}
\end{quote}

\begin{quote}
After the strawberries the child came out in a rash.
\end{quote}

Feminine.

\begin{cousinweb}
\emph{\ru{сыпь}} is built on \emph{\ru{сыпать}} 'to pour, to strew', which is inherited from Proto-Slavic \emph{sypati} (EtymDB): a rash is strewn over the skin.
\end{cousinweb}

\begin{grammarlens}[title={What you've built}]
One more word for telling, explaining and describing at B1.
\end{grammarlens}

\subsection*{Your turn}
\begin{itemize}
\raggedright
  \item \emph{Say it:} \emph{syp}
  \item \emph{Say it:} \emph{syp}, once more
  \item \emph{Say it:} a sentence of your own with \emph{syp}
  \item \emph{Recall:} say the Russian for to go hoarse, then the Russian for to tickle, then say \emph{syp} again
\end{itemize}

\begin{tcolorbox}[breakable,colback=teal!4,colframe=teal!35!black,title={Before you move on}]
What does \ru{сыпь} mean? (\textquotedblleft{}A rash\textquotedblright{}.) Say it once more.
\end{tcolorbox}
\section[\texorpdfstring{\ru{зуд} (zud)}{зуд (zud)}]{\ru{зуд} (zud) — an itch}

\label{lesson:RU-C487-zud}



\begin{quote}
Before the new one: say the Russian for to tickle, then the Russian for a rash.
\end{quote}

\subsection*{You'll want to know: \ru{зуд}}
\textbf{\ru{зуд}} (\emph{zud}) — \textquotedblleft{}an itch\textquotedblright{}.

\begin{quote}
\ru{Мазь} \ru{сняла} \ru{зуд}.
\end{quote}

\begin{quote}
\emph{Maz snyalá zud.}
\end{quote}

\begin{quote}
The ointment relieved the itch.
\end{quote}

Masculine.

\begin{cousinweb}
EtymDB lists \emph{\ru{зуд}} without a source, and the origin is not settled.
\end{cousinweb}

\begin{grammarlens}[title={What you've built}]
One more word for telling, explaining and describing at B1.
\end{grammarlens}

\subsection*{Your turn}
\begin{itemize}
\raggedright
  \item \emph{Say it:} \emph{zud}
  \item \emph{Say it:} \emph{zud}, once more
  \item \emph{Say it:} a sentence of your own with \emph{zud}
  \item \emph{Recall:} say the Russian for to tickle, then the Russian for a rash, then say \emph{zud} again
\end{itemize}

\begin{tcolorbox}[breakable,colback=teal!4,colframe=teal!35!black,title={Before you move on}]
What does \ru{зуд} mean? (\textquotedblleft{}An itch\textquotedblright{}.) Say it once more.
\end{tcolorbox}
\section[\texorpdfstring{\ru{чесаться} (chesátsya)}{чесаться (chesátsya)}]{\ru{чесаться} (chesátsya) — to itch; to scratch oneself}

\label{lesson:RU-C487-chesatsya}



\begin{quote}
Before the new one: say the Russian for a rash, then the Russian for an itch.
\end{quote}

\subsection*{You'll want to know: \ru{чесаться}}
\textbf{\ru{чесаться}} (\emph{chesátsya}) — \textquotedblleft{}to itch; to scratch oneself\textquotedblright{}.

\begin{quote}
\ru{После} \ru{укуса} \ru{комара} \ru{рука} \ru{очень} \ru{чешется}.
\end{quote}

\begin{quote}
\emph{Pósle ukúsa komará ruká óchen chéshetsya.}
\end{quote}

\begin{quote}
After the mosquito bite my hand itches a lot.
\end{quote}

Imperfective; present \ru{чешусь}, \ru{чешешься} (\ru{с} becomes \ru{ш}). Perfective \ru{почесаться}.

\begin{cousinweb}
\emph{\ru{чесаться}} is the reflexive of \emph{\ru{чесать}} 'to scratch, to comb', which is inherited from Proto-Slavic \emph{česati} (EtymDB).
\end{cousinweb}

\begin{grammarlens}[title={What you've built}]
One more word for telling, explaining and describing at B1.
\end{grammarlens}

\subsection*{Your turn}
\begin{itemize}
\raggedright
  \item \emph{Say it:} \emph{chesátsya}
  \item \emph{Say it:} \emph{chesátsya}, once more
  \item \emph{Say it:} a sentence of your own with \emph{chesátsya}
  \item \emph{Recall:} say the Russian for a rash, then the Russian for an itch, then say \emph{chesátsya} again
\end{itemize}

\begin{tcolorbox}[breakable,colback=teal!4,colframe=teal!35!black,title={Before you move on}]
What does \ru{чесаться} mean? (\textquotedblleft{}To itch; to scratch oneself\textquotedblright{}.) Say it once more.
\end{tcolorbox}
\section[\texorpdfstring{\ru{прививка} (privívka)}{прививка (privívka)}]{\ru{прививка} (privívka) — a vaccination, a jab}

\label{lesson:RU-C487-privivka}



\begin{quote}
Before the new one: say the Russian for an itch, then the Russian for to itch.
\end{quote}

\subsection*{You'll want to know: \ru{прививка}}
\textbf{\ru{прививка}} (\emph{privívka}) — \textquotedblleft{}a vaccination, a jab\textquotedblright{}.

\begin{quote}
\ru{Перед} \ru{зимой} \ru{мы} \ru{сделали} \ru{прививку} \ru{от} \ru{гриппа}.
\end{quote}

\begin{quote}
\emph{Péred zimóy my sdélali privívku ot gríppa.}
\end{quote}

\begin{quote}
Before the winter we had a flu jab.
\end{quote}

Feminine. \emph{\ru{Сделать} \ru{прививку} \ru{от}} + genitive — to be vaccinated against.

\begin{cousinweb}
\emph{\ru{прививка}} is built on \emph{\ru{прививать}} 'to graft, to inoculate', from \emph{\ru{вить}} 'to twist, to weave', which is inherited from Proto-Slavic \emph{viti} (EtymDB).
\end{cousinweb}

\begin{grammarlens}[title={What you've built}]
One more word for telling, explaining and describing at B1.
\end{grammarlens}

\subsection*{Your turn}
\begin{itemize}
\raggedright
  \item \emph{Say it:} \emph{privívka}
  \item \emph{Say it:} \emph{privívka}, once more
  \item \emph{Say it:} a sentence of your own with \emph{privívka}
  \item \emph{Recall:} say the Russian for an itch, then the Russian for to itch, then say \emph{privívka} again
\end{itemize}

\begin{tcolorbox}[breakable,colback=teal!4,colframe=teal!35!black,title={Before you move on}]
What does \ru{прививка} mean? (\textquotedblleft{}A vaccination, a jab\textquotedblright{}.) Say it once more.
\end{tcolorbox}
\section[\texorpdfstring{\ru{заражаться} (zarazhátsya)}{заражаться (zarazhátsya)}]{\ru{заражаться} (zarazhátsya) — to get infected}

\label{lesson:RU-C487-zarazhatsya}



\begin{quote}
Before the new one: say the Russian for to itch, then the Russian for a vaccination.
\end{quote}

\subsection*{You'll want to know: \ru{заражаться}}
\textbf{\ru{заражаться}} (\emph{zarazhátsya}) — \textquotedblleft{}to get infected\textquotedblright{}.

\begin{quote}
\ru{Дети} \ru{быстро} \ru{заражаются} \ru{друг} \ru{от} \ru{друга}.
\end{quote}

\begin{quote}
\emph{Déti býstro zarazháyutsya drug ot drúga.}
\end{quote}

\begin{quote}
Children catch things from each other very quickly.
\end{quote}

Imperfective; perfective \ru{заразиться} (\ru{заражусь}). \emph{\ru{Заражаться} \ru{от} \ru{кого}} — to catch it from someone.

\begin{cousinweb}
\emph{\ru{заражаться}} is built on \emph{\ru{зараза}} 'infection'; EtymDB lists \emph{\ru{зараза}} and \emph{\ru{заразить}} without a source, and they are usually traced to \emph{\ru{разить}} 'to strike'.
\end{cousinweb}

\begin{grammarlens}[title={What you've built}]
One more word for telling, explaining and describing at B1.
\end{grammarlens}

\subsection*{Your turn}
\begin{itemize}
\raggedright
  \item \emph{Say it:} \emph{zarazhátsya}
  \item \emph{Say it:} \emph{zarazhátsya}, once more
  \item \emph{Say it:} a sentence of your own with \emph{zarazhátsya}
  \item \emph{Recall:} say the Russian for to itch, then the Russian for a vaccination, then say \emph{zarazhátsya} again
\end{itemize}

\begin{tcolorbox}[breakable,colback=teal!4,colframe=teal!35!black,title={Before you move on}]
What does \ru{заражаться} mean? (\textquotedblleft{}To get infected\textquotedblright{}.) Say it once more.
\end{tcolorbox}
